\section{Discussion}

The empirical findings in the previous section can be read in two ways: as descriptive statistics about a dataset, or as a portrait of a community caught in the process of forming. This section takes the second reading.

\subsection{Early activity and community formation}

The late-2024 break visible in the corpus can be read not only as a change in volume, but also as part of the early formation of a Slovene Bluesky public during a broader period of platform change and migration. Early activity on newly adopted platforms often includes meta-platform orientation: posts about the platform itself, comparisons with older platforms, and negotiation of emerging features, norms, and expectations. Similar dynamics were also visible in public discussion around platforms such as Mastodon and Threads, where users likewise frequently compared them to Twitter/X; a comparable orientation is plausible here as well, even if it was not directly measured in the present study. From a user perspective, a platform arguably becomes fully adopted only once this phase fades into the background and ordinary use takes over: when users stop foregrounding how the platform works and simply use it. The late-2024 surge and the retention patterns observed in Section~\ref{sec:empirical} suggest that, for many authors, Bluesky was by that point becoming a more regular space of participation rather than remaining only a novelty. Whether this transition is also visible in topic choice, register, or interactional tone remains a question for future work.

\subsection{Dialogue over broadcasting}

The corpus is dominated by replies, and this pattern is not confined to a small number of marginal or highly active accounts. In that sense, Slovene Bluesky use appears to be oriented more towards exchange than towards one-way broadcasting. Because reply-dominant use is already visible before the late-2024 surge, the post-surge expansion is better understood as a continuation and strengthening of an existing interactional tendency than as a shift in kind.

\subsection{Small but concentrated}

The Slovene Bluesky space captured here is small and highly concentrated, with a limited core of authors contributing a large share of the corpus. Its communicative profile is shaped less by broad hashtag coordination than by conversational exchange, recurring news-media linking, and widely shared informal practices such as GIF use. Taken together, these patterns suggest a small public oriented more around interaction and repeated content-sharing than around large-scale topic publics. At the same time, the size of the corpus limits stronger claims about internal subgroup structure.

\subsection{ATProto as a research infrastructure}

From a methodological standpoint, the experience of building this corpus confirms that ATProto’s open design has real consequences for research. Public synchronisation and repository mechanisms made a backfill of the discovered Slovene-linked authors achievable without special research access \parencite{atproto-sync,atproto-streaming}. For researchers working on small or under-resourced languages, where corpora are scarce and institutional resources are limited, this openness matters. The Slovene community on Bluesky may be small, but the corpus captures much of the community’s publicly available posting history.
